\section{Introduction}

This report documents the final stage of a project whose goal is to turn a hiring
need, expressed in natural language, into a coherent pack of HR documents.
Mini-Assignment~1 adapted a small, open-source language model to the domain of IT job
postings. Mini-Assignment~2 taught that checkpoint to follow recruiter instructions
through supervised fine-tuning and DPO-based preference alignment. The present work,
\textit{job2cool}, integrates the resulting aligned model into a complete application.
Its user-facing persona is Diana, an HR assistant that interacts with the recruiter to
produce a Job Offer, a set of Technical Interviews, an Onboarding Plan, and a Cultural
and Team Fit assessment, all grounded in a company knowledge base, written live into
editable documents, and cited with click-through links to the exact source PDF page and
bounding box.

At the core of the system sits \verb|ma2-360m-dpo-b01|, the SmolLM2-360M checkpoint
adapted in Mini-Assignment~1 and DPO-aligned in Mini-Assignment~2. Around it, job2cool
adds a general-purpose orchestrator, \verb|gemma-4| (E4B-class, 128\,K context), and four
additional course components beyond the alignment work already done, two more than the
assignment requires:

\begin{enumerate}
  \item \textbf{Retrieval-augmented generation}: concurrent vector and graph search
    over a knowledge base of ten role-specific onboarding domains and a 210,048-CV
    candidate corpus, grounding output in company material and enabling clickable
    citations to source PDFs with page-level bounding-box overlays.
  \item \textbf{Tool use and agentic orchestration}: \verb|gemma-4| LLM-decides each
    pipeline step (from conversational memory resolution, through section
    classification, to per-section retrieval and composition) with no rule- or
    regex-based routing.
  \item \textbf{LLM-as-a-judge}: every response is scored for faithfulness and
    answer-relevance against the retrieved evidence, surfaced per turn in the UI.
  \item \textbf{Performance optimisation}: llama.cpp quantised inference under a model
    router, plus a root-caused configuration fix for a multi-slot KV-buffer memory
    failure that had crashed the host four times during development.
\end{enumerate}

A fifth course technique, advanced reasoning, is present but counted as supporting
rather than headline: \verb|gemma-4| emits inline chain-of-thought that the orchestrator
both consumes (for its routing decisions) and surfaces live in a Thinking panel
(Section~\ref{sec:cot}). All models in the request path are open-source and quantised; no
closed API is used anywhere in the serving path.

The system is functional end-to-end, with two clear entry points: a recruiter-facing web
UI and a command-line candidate-ingestion tool. In the UI, Diana identifies the role from
conversational context, resolves the matching knowledge-base domain, determines which of
the four documents were requested, retrieves evidence concurrently from vector and graph
indices, streams a reasoning-aware introduction, and generates each document
progressively in a tabbed workspace whose citations open the source PDF at the cited page
with a bounding-box highlight. The interface is a left-nav-rail application (Workspace,
Projects, Candidates, Agents, Skills, Tools, and Knowledge Base) served over fifteen
request-path containers across three Docker networks, authenticated via Google identity
through an \verb|oauth2-proxy| gateway. A Candidates browser gives direct access to the
indexed CV corpus; a Projects view provides private and shared project management with
full-fidelity replay of every past turn, including its reasoning and grounding panels;
and a live infrastructure health map monitors all services event-driven.

Although development and evaluation focused on IT recruitment, the architecture does not
depend on characteristics specific to this sector. The domain-organised knowledge base,
the domain-resolution mechanism, and the orchestration pipeline are all independent of
the concrete content of those domains, so the same system could in principle be adapted
to other recruitment areas by replacing the knowledge base and the specialised drafting
model.
